The scratch-wound healing assay is a cornerstone of \textit{in vitro} cell migration studies, yet automated quantification remains limited by single-frame independence, training-data requirements, and the absence of biological constraints in existing tools. Here we present \textbf{Quantification}, a fully automated computational pipeline that segments wound boundaries from phase-contrast time-lapse microscopy using local variance texture analysis and enforces a biologically motivated monotonic closure constraint across the temporal sequence. Unlike deep-learning approaches, Quantification requires no annotated training data and operates on standard phase-contrast images in $O(HW)$ time per frame. The pipeline comprises a seven-stage single-frame detector---local variance mapping, Y-band detection, column-wise edge extraction, RANSAC-based outlier rejection, Savitzky--Golay smoothing, vectorized mask reconstruction, and geometric validation---followed by a stateful temporal integrator available in two modes: a per-column Kalman filter that fuses detector observations with monotonic predictions weighted by per-frame confidence, and a deterministic hard constraint. We validated Quantification on 325 baseline images from two independent scratch-wound experiments in immortalized skeletal muscle cell lines (iMC ISOR544C and iMC MUTR544C) treated with an ALK5 inhibitor (TGF-$\beta$ receptor~I antagonist) or vehicle control (DMSO). Expert manual review confirmed 98.9\% detection accuracy on 285 wound-bearing images, with a mean Dice coefficient of 0.91 against hand-drawn ground truth ($n = 5$). The automatic quality-control algorithm achieved perfect sensitivity (1.00) for wound-bearing images. Downstream analysis of 38 validated trajectories at 0.1~mM concentration revealed no statistically significant differences in wound closure speed or edge displacement between ALK5i-treated and DMSO-treated samples after Bonferroni correction, though effect sizes were consistently negative (Cohen's $d = -0.26$ to $-0.43$), suggesting a trend toward reduced migration that may become detectable with larger sample sizes or longer observation windows. No condition reached 50\% closure within the 24-hour observation period. The pipeline processes a typical 20-frame trajectory in under 10~seconds on commodity hardware and is released as open-source Python code. Quantification addresses a critical methodological gap by combining training-free segmentation with physics-informed temporal constraints, producing smooth, strictly non-increasing area trajectories that are directly interpretable without post-hoc smoothing.
